\section{The original homogeneous line coefficient has the original analytic value}

\begin{lemma}[Whole coordinate-polynomial identification of original homogeneous line coefficients]
\label{lem:fc-original-homogeneous-section-coordinate}
\leanok
\lean{EnvelopeProperty.ProjBerk.homogeneousSection_coefficient_eq_coordinate}
\uses{lem:fc-original-homogeneous-form-coefficients}
For every original homogeneous form $F$ of degree $d$ over an arbitrary
coefficient ring, its whole original line-frame coefficient $F/X_j^d$ is
exactly the actual original affine polynomial section associated to the
original dehomogenization of $F$ at $X_j=1$.
\end{lemma}
\begin{proof}
The actual homogeneous-chart dehomogenization sends the genuine homogeneous
fraction to the original coordinate polynomial.  The two whole chart ring
maps are genuine inverses.  Apply the actual homogeneous away-to-section
map to this inverse identity.  The actual native coordinate-section
isomorphism is its composite with exactly the same chart ring inverse,
which gives the literal original whole section equality.
\end{proof}

\begin{lemma}[Every value-one chart preserves the original homogeneous form value]
\label{lem:fc-original-homogeneous-section-value}
\leanok
\lean{EnvelopeProperty.Berkovich.RelativeNormalizedProjectiveSpectrum.nativeHomeomorph_eval_homogeneousSection_coefficient}
\uses{lem:fc-original-homogeneous-section-coordinate}
For an actual prescribed-base normalized projective seminorm and every
coordinate $X_j$ of value one, evaluation of the actual original homogeneous
line coefficient on its native valued scheme point is exactly the original
full homogeneous seminorm value of $F$.  The coefficient ring and prescribed
base seminorm are arbitrary; this equality does not require a bound on the
coefficient ring.
\end{lemma}
\begin{proof}
Replace the whole original coefficient by the actual coordinate-polynomial
section using the preceding equality.  The true native projective chart
comparison preserves every original affine polynomial value.  Original
homogeneous recomposition gives $F=X_j^d F|_{X_j=1}$ in the actual localization;
since $X_j$ has value one, the original two polynomial values coincide.
\end{proof}
